\documentclass{article}

\usepackage{amsmath}

\title{Ch. 31 HW \#2 Solutions}
\author{Principles of Macroeconomics}
\date{\today}

\begin{document}

\maketitle

\section*{Quick Quiz}

\begin{enumerate}

\setcounter{enumi}{1}
\item (d) -- We can find the price level using nominal and real GDP.
	\begin{gather*}
	PY = \text{GDP}_{nom} \\
	200P = 400 \\
	P = 2
	\end{gather*}
	Then we can find velocity using the quantity equation.
	\begin{gather*}
	MV = PY \\
	100V = 400 \\
	V = 4
	\end{gather*}

\item (b) -- A major assumption of the quantity theory is that velocity is stable. The money supply, price level, and output all change over time

\item (b) -- Hyper inflations typically occur when the government spends a bunch of money and pays for it by "printing" money, a.k.a. expanding the money supply.

\item (a) -- If the central bank increases the rate of money growth, then the quantity theory of money says that inflation will increase. If inflation increases, the Fisher effect says that nominal interest rates will increase.

\end{enumerate}

\section*{Questions for Review}

\begin{enumerate}

\setcounter{enumi}{3}
\item Inflation is like a tax in two ways. First, inflation can help fund a government because a government can pay for its spending by printing new currency and inflating the price level. Second, the cost of inflation is paid by the public because it erodes the value of their money. 

\item In the long term, an increase in the inflation rate causes an increase in the expected inflation rate. According to the Fisher effect, that increases the nominal interest rate approximately one-to-one. It doesn't affect the real interest rate because of monetary neutrality.

\end{enumerate}

\section*{Problems and Applications}

\begin{enumerate}

\item

	\begin{enumerate}
	
	\item Price level:
		\begin{gather*}
		\text{GDP}_{Nom} = PY \\
		10 = 5P \\
		\$2 \text{ trillion} = P
		\end{gather*}
		Velocity:
		\begin{gather*}
		MV = PY \\
		0.5V = 10 \\
		V = 20
		\end{gather*}
		
	\item The price level will decrease by 5\%.
		\begin{gather*}
		MV = PY \\
		0.5(20) = P(1.05)(5) \\
		\frac{2}{1.05} = P \\
		1.9 = P \\
		\% \Delta P = \frac{1.9 - 2}{2} = -5\%
		\end{gather*}
		
	\item It should increase the money supply by 5\%, so $M = 1.05(0.5) = 0.525$. That's \$525 billion.
	
	\item It should increase the money supply to \$577.5 billion. 
		\begin{gather*}
		MV = PY \\
		20M = 1.1(2)(1.05)(5) \\
		M = \frac{11.55}{20} \\
		M = 0.5775
		\end{gather*}
		
	\end{enumerate}

\end{enumerate}

\end{document}